\documentclass[11pt,a4paper]{article}
\usepackage[utf8]{inputenc}
\usepackage[T1]{fontenc}
\usepackage{lmodern}
\usepackage{geometry}
\geometry{margin=24mm}
\usepackage{microtype}
\usepackage{booktabs}
\usepackage{tabularx}
\usepackage{array}
\usepackage{enumitem}
\usepackage{longtable}
\usepackage{xcolor}
\usepackage{float}
\usepackage{hyperref}
\hypersetup{colorlinks=true,linkcolor=black,citecolor=blue,urlcolor=blue}
\newcolumntype{Y}{>{\raggedright\arraybackslash}X}
\newcommand{\LCDM}{$\Lambda$CDM}
\newcommand{\Ez}{$E(z)$}
\newcommand{\Hz}{$H(z)$}
\newcommand{\HzEq}{$H(z)=H_0E(z)$}
\title{A Claim-Bounded Failure-Boundary Audit of an A10-Derived Hubble-Tension Scaffold\\\large From Candidate Proxy to Physical-Bridge Non-Establishment}
\author{Keiji Yoshimura\\ \small Independent Researcher}
\date{Version v0.1.0-claim-bounded-failure-boundary-audit\\June 2026}
\begin{document}
\maketitle

\begin{abstract}
This manuscript reports a claim-bounded failure-boundary audit of an A10-derived Hubble-tension scaffold. The paper does not claim a solution to the Hubble tension, a replacement for \LCDM, an observational discovery, or a real-data likelihood result. The audited chain was designed to determine whether A10-origin materials could be reduced to a form suitable for later cosmological comparison. The terminal result is negative but informative: a candidate proxy, an internal diagnostic curve, a contract-only bridge scaffold, and a bounded hypothesis scaffold were retained; however, no first-principles physical background-expansion equation \Ez{} or \Hz{} was established inside the audited artifacts. Consequently, the chain did not proceed to \LCDM{} reference smoke, real-data likelihood testing, matched baseline comparison, MCMC, posterior comparison, or evidence evaluation. The principal contribution is therefore not a new cosmological model, but a documented failure-localization record: the unresolved boundary lies between candidate scaffold material and a physically typed background cosmology interface. This record is preserved as a methodological closeout and as input to a possible future rebuild of A10 Core as a structured-prior and strict-audit research operating procedure rather than as a validated physical theory.
\end{abstract}

\noindent\textbf{Keywords:} Hubble tension; failure-boundary audit; claim boundary; A10 methodology; cosmology scaffold; negative result; research governance.

\section{Introduction}

The Hubble tension motivates many proposals that attempt to modify the early or late expansion history of the Universe. Within the author's A10 project, a Hubble-tension branch was previously explored as a structured candidate involving a localized early-time transfer, stability-gated selectivity language, and a high-$H_0$ target region. Earlier A10 Core documents deliberately limited the public claim: A10 Core was framed as a structured-prior methodology, not as a universal physical law or a domain-validated cosmological theory \cite{YoshimuraA10Core2026}. The same documents also stated that stronger cosmological claims would require independent reproduction, likelihood comparison, robustness testing, and expert review.

This paper records the subsequent audit rather than promoting the earlier scaffold. The audit asked a stricter question: after deconstruction, could the A10-derived candidate proxy be connected to a first-principles physical background expansion equation, such as \Ez{} or \HzEq, strongly enough to justify a later \LCDM{} reference smoke or real-data likelihood comparison? The answer in the audited artifacts was no. This is not a numerical defeat by data, because the audit stopped before real data and before likelihood evaluation. It is a structural negative result: the physical bridge required to enter cosmological model comparison was not established.

The purpose of this manuscript is to make that stop point explicit. It records what survived, what did not, and why the chain stopped. It also seals the public claim boundary so that the archive cannot be misread as a Hubble-tension solution or as an observational model-selection result.

\section{A10 Core context and audit object}

The relevant A10 Core position is methodological. A10 Core v0.2.1 describes the framework as a structured-prior and claim-boundary methodology for constructing diagnostic branches, not as a single equation that overrides domain science \cite{YoshimuraA10Core2026}. Its operational sentence is: bold hypothesis, strict process, bounded claim. Under that rule, unusual hypotheses may enter an exploratory workflow, but only evidence-bounded claims may exit it.

The object audited here was not A10 Core as a general workflow. The audited object was narrower: an A10-derived Hubble-tension scaffold, including candidate proxy material, bridge-hypothesis attempts, and terminal handoff packages generated after a deconstruction workflow. The pre-audit seed dossier had already begun stripping A10-specific identity from the H0 branch, separating carry-forward candidates from deferred, ablated, or archived components \cite{SeedDossier2026}. Table~\ref{tab:strip} summarizes that stripping ledger.

\begin{table}[H]
\centering
\small
\caption{A10 element stripping ledger used as pre-audit source material.}
\label{tab:strip}
\begin{tabularx}{\textwidth}{p{16mm}p{31mm}Y}
\toprule
ID & Status & Component \\
\midrule
C01 & Carry forward & Generic localized early-time transfer kernel near $z \approx 5000$. \\
C02 & Carry forward & Quasi-vacuum to stiff-fluid reduced surrogate. \\
C03 & Carry forward & Perturbation-level validation requirement. \\
C04 & Carry forward & Residual-control objective under \LCDM{} and benchmark comparators. \\
C05 & Defer or ablate & Generic stability-gated activation; fixed $\chi \approx 137$ only as ablation coordinate. \\
C06 & Defer or ablate & Spectral-Horndeski stability-trigger idea. \\
C07 & Archive only & Root-geometric / E8 / ADE selectivity language. \\
C08 & Defer or ablate & Late-time drag / $S_8$ companion sector. \\
C09 & Defer or ablate & High-$H_0$ candidate behavior. \\
\bottomrule
\end{tabularx}
\end{table}

The audit therefore did not attempt to preserve A10 branding. It attempted to determine whether any neutral residue could become a physically typed cosmological seed. The answer was bounded: some scaffold material remained, but the physically required bridge remained absent.

\section{Audit protocol}

The protocol deliberately separated allowed documentation operations from forbidden scientific operations. Allowed operations included hash verification, extraction of candidate specifications, symbolic guard checks on toy grids, derivation-gap matrices, no-go gate records, and terminal handoff synthesis. Forbidden operations included real-data ingestion, parser or adapter execution for external datasets, likelihood or chi-square evaluation, matched \LCDM{} baseline execution, MCMC, posterior comparison, evidence metrics, and public solution claims.

This design was intentional. If a physical \Ez{}/\Hz{} bridge could not be established, then moving to real-data likelihood or \LCDM{} comparison would only create a false appearance of rigor. The audit therefore used \texttt{D2\_READY=False} as the decisive gate condition.

\section{Deconstruction chain summary}

Table~\ref{tab:chain} summarizes the compact terminal chain reported in this repository. Earlier fine-grained modes were consolidated into compound packages when their outcomes were expected documentation or guard outcomes. Unexpected results would have required separate re-inspection; none of the terminal packages authorized real-data comparison.

\begin{table}[H]
\centering
\small
\caption{Audit-chain summary.}
\label{tab:chain}
\begin{tabularx}{\textwidth}{p{25mm}p{39mm}Yp{18mm}}
\toprule
Package & Role & Main result & D2 ready \\
\midrule
Phase 1.5 & A10 stripping and neutral seed dossier & A10-origin components were split into carry-forward, defer/ablate, and archive-only categories. & No \\
F2 & Symbolic-contract guard smoke & A normalized symbolic scaffold passed finite/positive toy-grid guards only. & False \\
F3 & Physical-status / derivation-gap matrix & The symbolic scaffold was separated from physical cosmology status. & False \\
F4--F7 & Compound derivation-gap and D2 gate & D2 prerequisites remained unmet; no likelihood or comparison was executed. & False \\
G0--G3 & First-principles-or-null decision & No first-principles physical \Ez{}/\Hz{} derivation was found in the current artifacts. & False \\
H0--H3 & Negative derivation closeout & Candidate scaffold was retained only as a bounded hypothesis scaffold. & False \\
I0--I3 & Human-readable terminal dossier & Expert handoff and no-go gate were synthesized. & False \\
J0--J3 & Terminal QA and public boundary & Public boundary was sealed; recommended next action was stop, restart, or expert handoff. & False \\
\bottomrule
\end{tabularx}
\end{table}

The key feature of this chain is the difference between a scaffold that is finite under toy guards and a physical cosmological model. The former is not sufficient for the latter. Passing a symbolic or contract-only guard means that the expression was well-behaved under a limited synthetic domain. It does not imply that the expression has a physical derivation, unit-consistent interpretation, observable-distance mapping, or likelihood-ready implementation.

\section{Main result: physical bridge non-establishment}

The terminal result can be stated in one sentence:

\begin{quote}
The A10-derived H0 scaffold retained bounded diagnostic material, but no first-principles physical \Ez{}/\Hz{} bridge was established in the audited artifacts.
\end{quote}

The main failure boundary is therefore not simply ``the theory failed'' in a vague sense. It is located at a specific interface:

\begin{center}
\begin{tabular}{c}
Candidate proxy / bounded scaffold \\
$\downarrow$ \\
First-principles physical \Ez{}/\Hz{} bridge equation \\
$\downarrow$ \\
Observable distance / likelihood / \LCDM{} comparison
\end{tabular}
\end{center}

The first line retained material. The second line was not established. The third line was therefore not entered. This is a structural negative result before observational comparison.

\section{Retained material}

The audit retained the following materials as bounded handoff artifacts:

\begin{itemize}[leftmargin=7mm]
\item a candidate proxy;
\item an internal diagnostic curve;
\item a contract-only bridge scaffold;
\item a distance-like integral guard, explicitly not an observable distance;
\item bridge-hypothesis sandbox candidates;
\item a failure-localization record;
\item a human-readable expert handoff synthesis.
\end{itemize}

These items are potentially useful as diagnostic records or as inputs to a future model-building discussion. They are not sufficient to define a cosmological model.

\section{Material not established}

The following items were not established:

\begin{itemize}[leftmargin=7mm]
\item a first-principles physical \Ez{} equation;
\item a first-principles physical \HzEq{} equation;
\item an observable cosmological distance model;
\item a SN/BAO/CMB likelihood;
\item a matched \LCDM{} baseline comparison;
\item MCMC, posterior comparison, or evidence evaluation;
\item a Hubble-tension solution claim;
\item a \LCDM{} replacement claim.
\end{itemize}

This distinction is the central claim boundary of the repository.

\section{Unauthorized operations remained zero}

The terminal audit chain maintained the intended firewall. The forbidden operations listed in Table~\ref{tab:zero} were not executed. This does not make the scaffold physically successful; it makes the negative result clean. The archive does not mix a scaffold-level stop with unauthorized observational claims.

\begin{table}[H]
\centering
\small
\caption{Forbidden operation counters at terminal closeout.}
\label{tab:zero}
\begin{tabularx}{\textwidth}{Yp{18mm}}
\toprule
Operation & Count \\
\midrule
Real-data use or ingestion & 0 \\
Parser / adapter execution for external data & 0 \\
Real likelihood or chi-square evaluation & 0 \\
Matched \LCDM{} baseline execution & 0 \\
MCMC execution & 0 \\
Posterior comparison & 0 \\
Evidence metric, Bayes factor, AIC, or BIC emission & 0 \\
Hubble-tension solution claim emission & 0 \\
\LCDM{} replacement claim emission & 0 \\
\bottomrule
\end{tabularx}
\end{table}

\section{Public claim boundary}

The public claim supported by this archive is narrow:

\begin{quote}
An A10-derived H0 scaffold was deconstructed and audited. Candidate proxy/scaffold material survived as bounded hypothesis material, but the physical \Ez{}/\Hz{} bridge required for cosmological comparison was not established. Therefore no \LCDM{} comparison, real-data likelihood, posterior, or evidence claim is made.
\end{quote}

The archive does not support statements that A10 solved the Hubble tension, defeated \LCDM, replaced \LCDM, or obtained observational support from SN/BAO/CMB likelihoods. It also does not support the opposite overclaim that \LCDM{} numerically defeated A10 in a matched likelihood contest, because such a contest was not executed.

\section{Implications for A10 Core rebuild}

The audit has an implication for A10 Core itself. If A10 Core is treated as a physical theory whose historical crown is the Hubble-tension branch, then this audit is damaging: the branch did not close as a physical background cosmology model. If A10 Core is instead treated as a strict research operating procedure, the result is more constructive. The workflow successfully stopped claim inflation, preserved the failure boundary, and demoted a tempting scaffold to bounded handoff material.

The rebuild implication is therefore clear. A10 Core should not use the cosmology branch as a success emblem. It may use the branch as a historical case study showing why strict audit and no-go ledgers are needed. In that sense, the cosmology branch becomes a record of self-audit rather than a crown.

Future A10-derived cosmology work, if any, should not begin by defending A10 identity. It should begin from a physically typed question: can any neutral residue be formulated as an effective-fluid deformation, modified background equation, or explicitly phenomenological diagnostic parameterization with unit consistency, observable mapping, and comparator discipline? If the answer is no, the cosmology route should remain closed.


\section{Why this is not a data-level rejection}

A common misunderstanding of this archive would be to read the result as if the A10-derived scaffold had been tested against data and had lost. That is not what happened. A data-level rejection would require a clearly defined physical model, a consistent mapping to observables, a matched baseline, a likelihood contract, and a comparison protocol. The present audit did not reach that level because the physical background interface was not established.

The negative result is therefore upstream of observational cosmology. It says that the audited artifacts did not contain enough physically typed structure to justify a comparison. In practical terms, this is a more basic stop than an unfavorable Bayes factor or a poor chi-square. It is the absence of a model-ready background equation.

This distinction matters for both humility and accuracy. The scaffold should not be advertised as having survived data, but neither should it be described as having been defeated by a data vector that was never loaded. Its status is: structurally incomplete for model comparison.

\section{Relation to a possible neutral H0 seed}

The Phase 1.5 seed dossier proposed a neutralized future frame: a compact pre-recombination transfer kernel evaluated without A10 identity, without a privileged $\chi \approx 137$ assumption, and without high-$H_0$ alone as a success criterion \cite{SeedDossier2026}. The current audit strengthens that instruction. If any future project restarts, it should not be called an A10 Hubble solution project. It should begin as a comparator-first residual-response audit.

A viable future seed would need at least four properties:

\begin{enumerate}[leftmargin=7mm]
\item a unit-consistent physical or phenomenological definition of the background expansion modification;
\item an explicit mapping from that modification to observable distances or perturbative observables;
\item benchmarks against \LCDM, $\Delta N_{\rm eff}$-like, early-dark-energy-like, no-gate, and generic-gate controls;
\item a stop rule that demotes the seed if it cannot pass perturbative, residual-cost, and prior-sensitivity tests.
\end{enumerate}

The seed dossier already recommended that root-geometric, E8, and ADE selectivity language be archived rather than carried forward as physical cosmological evidence. The present audit makes that recommendation even stronger. Such language may remain historically useful, but it should not be part of any future physically typed H0 model unless independently rebuilt by domain specialists.

\section{Decision matrix for next actions}

Table~\ref{tab:next} summarizes the appropriate next-action matrix. The table is included to prevent a common failure mode: treating a negative bridge result as an invitation to continue directly into a more complex pipeline. The audit recommends the opposite. Complexity should not be added until the missing interface is addressed.

\begin{table}[H]
\centering
\small
\caption{Next-action decision matrix.}
\label{tab:next}
\begin{tabularx}{\textwidth}{p{28mm}Yp{29mm}}
\toprule
Action & Condition & Status \\
\midrule
Stop and archive & No physical \Ez{}/\Hz{} bridge exists in the audited artifacts. & Recommended and sufficient. \\
Expert handoff & A reviewer is asked whether any scaffold can be reformulated as an effective-fluid, modified-background, or phenomenological model. & Allowed with boundary. \\
A10 Core rebuild & The failure log is used to demote cosmology from success emblem to self-audit case study. & Recommended. \\
Neutral H0 seed project & A new typed equation or phenomenological parameterization is constructed independently of A10 branding. & Possible but not started here. \\
D2 / \LCDM{} reference smoke & A physical background interface, observable mapping, and comparator contract are established. & Not allowed yet. \\
Public solution claim & Any claim that the Hubble tension is solved or that \LCDM{} is replaced. & Forbidden. \\
\bottomrule
\end{tabularx}
\end{table}

\appendix
\section{Terminal decision matrix}

The terminal J0--J3 package preserved the following matrix in human-readable form \cite{J0J3Audit2026}. The table is reproduced here as a compact publication summary.

\begin{table}[H]
\centering
\small
\caption{Terminal decision matrix reproduced from the terminal closeout package.}
\begin{tabularx}{\textwidth}{p{43mm}p{24mm}Y}
\toprule
Item & Decision & Comment \\
\midrule
Candidate proxy & Retain & Internal diagnostic / bounded scaffold only. \\
Contract-only bridge scaffold & Retain & Not a physical equation. \\
Distance-like integral guard & Retain & Not an observable cosmological distance. \\
First-principles physical \Ez{} derivation & Not found & Not established inside current artifacts. \\
Physical \HzEq{} equation & Not established & Unit, derivation, and observable mapping remain unresolved. \\
D2 \LCDM{} reference smoke & No-go & Preconditions unmet. \\
Real-data likelihood & No-go & Remained forbidden. \\
MCMC / posterior / evidence & No-go & Remained forbidden. \\
Public solution claim & Forbidden & Would violate the claim boundary. \\
Expert handoff & Allowed & Only as bounded negative/diagnostic handoff. \\
\bottomrule
\end{tabularx}
\end{table}

\section{No-go statement for reuse}

For repository reuse, the shortest safe statement is:

\begin{quote}
This archive documents a claim-bounded failure-boundary audit. It retained candidate scaffold material but did not establish a first-principles physical \Ez{}/\Hz{} bridge. Therefore it does not report a real-data likelihood, \LCDM{} comparison, MCMC posterior, evidence metric, Hubble-tension solution, or \LCDM{} replacement claim.
\end{quote}

Any derivative document should preserve this statement or an equivalent boundary.


\section{Limitations}

This manuscript is limited by design. It does not reproduce the full historical A10 cosmology chain, it does not load external data, and it does not independently run CLASS, Cobaya, or any full CMB/BAO/SN likelihood. The audit is therefore not an observational cosmology paper. It is a methodological and archival closeout paper. Its value lies in exact claim separation: scaffold retained; physical bridge not established; comparison not authorized.

\section{Conclusion}

The A10-derived H0 scaffold did not mature into a likelihood-ready cosmological model in the audited chain. The central missing piece was a first-principles physical \Ez{}/\Hz{} bridge equation. Because that bridge was not established, the audit stopped before \LCDM{} reference smoke, real-data likelihood, MCMC, posterior comparison, or evidence evaluation.

This is a negative result, but not an empty one. The archive preserves the exact failure boundary and the remaining bounded materials. It also demonstrates the stricter A10 posture in practice: bold hypotheses may enter, but unsupported claims do not exit. The correct final status is therefore not victory and not observational defeat, but claim-bounded failure-localization.


\section*{Repository and Data Availability}

The repository package associated with this failure-boundary audit is available at:

\begin{center}
\url{https://github.com/yokken0907/a10-derived-h0-scaffold-failure-boundary-audit}
\end{center}

The repository is provided as an audit archive and reproducibility aid. It is not an independent publication claim, and no repository release DOI or Zenodo DOI is assigned in this manuscript. The repository does not report a real-data likelihood, \LCDM{} comparison, MCMC posterior, evidence metric, Hubble-tension solution, or \LCDM{} replacement claim. Any future repository release or archival DOI should preserve the same claim boundary and should not be treated as a substitute for the present preprint record.

\section*{Acknowledgments}

The author conceptually designed the study and formulated the core audit questions. The author acknowledges the use of large language models, including ChatGPT, as computational and editorial assistants for manuscript organization, code-oriented reasoning, workflow packaging, and language refinement under the author's direct supervision and validation. The author remains responsible for the final content, claim boundary, and interpretation.

\begin{thebibliography}{11}
\bibitem{YoshimuraA10Core2026} K. Yoshimura, \emph{A10 Core Foundation v0.2.1: Theory-Freedom / Strict-Audit Clarification, and a Localized Early-Time Cosmology Case Study}, repository manuscript, May 2026.
\bibitem{SeedDossier2026} K. Yoshimura, \emph{A10 Core to New H0 Project Seed Dossier, Phase 1.5}, repository artifact, 2026.
\bibitem{J0J3Audit2026} K. Yoshimura, \emph{A10 Cosmology v0.5.0 Mode J0--J3 Compound Terminal QA, Public Boundary, and Restart Handoff}, repository artifact, 2026.
\bibitem{Planck2020} Planck Collaboration, N. Aghanim et al., ``Planck 2018 results. VI. Cosmological parameters,'' \emph{Astronomy and Astrophysics} 641, A6 (2020). DOI: 10.1051/0004-6361/201833910.
\bibitem{Riess2022} A. G. Riess et al., ``A Comprehensive Measurement of the Local Value of the Hubble Constant with 1 km s$^{-1}$ Mpc$^{-1}$ Uncertainty from the Hubble Space Telescope and the SH0ES Team,'' \emph{Astrophysical Journal Letters} 934, L7 (2022). DOI: 10.3847/2041-8213/ac5c5b.
\bibitem{Karwal2016} T. Karwal and M. Kamionkowski, ``Early dark energy, the Hubble-parameter tension, and the string axiverse,'' \emph{Physical Review D} 94, 103523 (2016). DOI: 10.1103/PhysRevD.94.103523.
\bibitem{Poulin2019} V. Poulin, T. L. Smith, T. Karwal, and M. Kamionkowski, ``Early Dark Energy Can Resolve the Hubble Tension,'' \emph{Physical Review Letters} 122, 221301 (2019). DOI: 10.1103/PhysRevLett.122.221301.
\bibitem{Hill2020} J. C. Hill, E. McDonough, M. W. Toomey, and S. Alexander, ``Early dark energy does not restore cosmological concordance,'' \emph{Physical Review D} 102, 043507 (2020). DOI: 10.1103/PhysRevD.102.043507.
\bibitem{CLASSII2011} D. Blas, J. Lesgourgues, and T. Tram, ``The Cosmic Linear Anisotropy Solving System (CLASS) II: Approximation schemes,'' \emph{Journal of Cosmology and Astroparticle Physics} 07, 034 (2011). DOI: 10.1088/1475-7516/2011/07/034.
\bibitem{Lesgourgues2011} J. Lesgourgues, ``The Cosmic Linear Anisotropy Solving System (CLASS) I: Overview,'' arXiv:1104.2932 [astro-ph.IM] (2011).
\bibitem{Alam2017} S. Alam et al., ``The clustering of galaxies in the completed SDSS-III Baryon Oscillation Spectroscopic Survey: cosmological analysis of the DR12 galaxy sample,'' \emph{Monthly Notices of the Royal Astronomical Society} 470, 2617--2652 (2017). DOI: 10.1093/mnras/stx721.
\end{thebibliography}

\end{document}
